%---
% title: "Why a low-SNR rank-one gradient moment does not justify EWC"
% date: 2026-06-18
% work_order: 8
% permalink: /posts/2026/06/rank-1-fisher-no-task-to-remember/
% tags:
%   - machine-learning
%   - continual-learning
%   - diffusion-models
%   - fisher-information
%---
\documentclass[11pt]{article}
\usepackage{publication-web}
\begin{document}

\href{https://arxiv.org/abs/2509.23593}{Avoid Catastrophic Forgetting with Rank-1 Fisher from Diffusion Models} forms the second moment \(F_{\mathrm E}=\mathbb E[gg^\top]\) of per-example diffusion-loss gradients and uses a rank-one approximation to \(F_{\mathrm E}\) in Elastic Weight Consolidation (EWC). The proposed justification has three separate problems: the linear proof contains an invalid matrix identity, the leading low-SNR term is independent of the old task, and \(F_{\mathrm E}\) is not generally Fisher information.

\section{What EWC must preserve}\label{what-ewc-must-preserve}

EWC adds a local old-task penalty to the new-task loss,

\[
L_B(\theta)+\frac{\lambda}{2}(\theta-\theta_A^\star)^\top F_A(\theta-\theta_A^\star).
\]

In its Bayesian derivation, \(F_A\) approximates the local curvature of the negative log posterior for task \(A\) near \(\theta_A^\star\) (\href{https://arxiv.org/abs/1612.00796}{Kirkpatrick et al.}). If \(F_A=\lambda_1vv^\top\), the penalty constrains only the scalar displacement \(v^\top(\theta-\theta_A^\star)\). Interpreting this penalty as old-task preservation therefore requires \(v\) to identify parameter changes to which task \(A\) is sensitive.

\section{The linear proof is incorrect}\label{the-linear-proof-is-incorrect}

The paper assumes a linear score model \(s_A(x)=Ax\) with target \(\gamma x\). Up to a positive scalar \(c\), the per-example loss and its gradient with respect to \(A\) are

\[
\ell(A;x)=\frac{c}{2}\lVert(A-\gamma I)x\rVert^2,
\qquad
\nabla_A\ell(A;x)=c(A-\gamma I)xx^\top.
\]

The proof replaces the last expression by \(c\lVert x\rVert^2(A-\gamma I)\), which would require \(xx^\top=\lVert x\rVert^2I\). This identity is false for \(d>1\). Except in one-dimensional or otherwise degenerate cases, the orientation of \(xx^\top\) varies across examples, so the vectorized gradients need not be collinear. The rank-one form of \(F_{\mathrm E}\) therefore does not follow even for the stated linear surrogate. Measurements on a nonlinear U-Net may still show a dominant empirical eigenvector, but they do not repair this derivation.

\section{Low SNR removes task identity}\label{low-snr-removes-task-identity}

For the variance-preserving forward process

\[
x_t=\sqrt{\bar\alpha_t}x_0+\sqrt{1-\bar\alpha_t}\epsilon,
\qquad \epsilon\sim\mathcal N(0,I),
\]

the noisy marginal converges to \(\mathcal N(0,I)\) as \(\bar\alpha_t\to0\). Its leading score and noise-prediction target are scaled copies of \(x_t\) whose coefficients depend on the noise schedule, not on the data distribution. This leading target contains no old-task identity, so the limiting argument does not establish that its gradient direction measures old-task sensitivity.

The implemented rank-one direction instead averages gradients over both data and sampled timesteps, \(\mathbb E_{x_0,t}[g(\theta;x_t,t)]\), and the continual-learning objective also includes generative distillation. Finite-SNR gradients and distillation can contain old-task information, but neither is identified by the task-independent limiting argument.

\section{A gradient second moment is not Fisher information}\label{gradient-moment-not-fisher}

The spectrum of \(F_{\mathrm E}=\mathbb E[\nabla\ell\nabla\ell^\top]\) does not in general equal the Fisher-information spectrum of a probabilistic model. For the conditional linear-Gaussian model

\[
x\sim\mathcal N(0,\Sigma_x),
\qquad
y\mid x,A\sim\mathcal N(Ax,\sigma^2I_m),
\]

the likelihood Fisher for \(A\) is

\[
F_{\mathrm{true}}=\frac{1}{\sigma^2}\Sigma_x\otimes I_m.
\]

If \(x\) is isotropic, this matrix is proportional to the identity rather than rank one. This calculation does not imply that the diffusion-loss gradient moment is useless; it shows that a dominant eigenvector of that moment does not establish a rank-one local KL geometry. \href{https://arxiv.org/abs/1905.12558}{Kunstner, Balles, and Hennig} analyze the distinction between gradient second moments and Fisher information more generally.

\section{What would establish the mechanism}\label{what-would-establish-the-mechanism}

A theoretical repair must retain the \(xx^\top\) factor or state conditions under which the resulting gradients are approximately collinear. Timestep-restricted EWC matrices, compared under the same replay and distillation procedure, would then show whether old-task preservation comes from low-, intermediate-, or high-SNR gradients. Until that link is established, the penalty should be described as empirically motivated by a gradient second moment, not derived from the stated low-SNR Fisher argument. \href{/files/improved-elastic-weight-consolidation.pdf}{IEWC} illustrates the distinct alternative of deriving an EWC penalty from preservation of earlier optimization updates.

\end{document}
